\subsection{予測可能補償過程の貼り合わせと jump bound}

各 $n$ について、$J^{\tau_n}$ の双対予測可能射影（補償過程）を $P^n$、martingale residualを
$Q^n=J^{\tau_n}-P^n$ と書く。局所構成で得た ordered overlap は
\begin{equation}
 m\leq n\quad\Longrightarrow\quad (P^n)^{\tau_m}\sim P^m
 \label{eq:compensator-ordered}
\end{equation}
である。一方、局所有限変動過程の貼り合わせが要求する形は
\begin{equation}
 (P^n)^{\tau_n\wedge\tau_m}\sim
 (P^m)^{\tau_n\wedge\tau_m}\qquad(n,m\in\mathbb N)
 \label{eq:compensator-minimum}
\end{equation}
である。式~\eqref{eq:compensator-ordered}で $m=n$ とすれば各 coordinate の自己停止が得られる。また一つの
full-measure set上で $(\tau_n)$ は増加するので、$m\leq n$ なら
$\tau_n\wedge\tau_m=\tau_m$ である。二重停止の吸収則と自然数の全順序を使えば
式~\eqref{eq:compensator-minimum}が従う。これにより ordered overlapを貼り合わせ定理のall-pairs入力へ移した。

式~\eqref{eq:compensator-minimum}から pointwise に eventually constantな局所極限を作ると、予測可能性は極限で保たれる。
compatibilityとscheduleのa.e. 増加性・exhaustionをそのまま用いることで、right continuityと
局所有限変動性がa.e. に得られる。そのpath-regularityの例外集合上で過程全体を零にすれば、
全sample pathでright-continuousかつ局所有限変動な予測可能version $P$ を得る。
compatibilityや $(\tau_n)$ 自体をpointwiseなdataへ置き換える必要はない。ただし各 $P^n_0=0$ は
a.e. の等式なので、貼り合わせだけでは $P_0=0$ も
a.e. にしか得られない。$\{P_0\neq0\}$ は時刻零で可測な零集合であり、そこでも過程全体を零にして
$P_0=0$ をpathwiseに正規化する。この二度目の正規化はpredictability、right continuity、
局所有限変動性、および各coordinateとのstopped indistinguishabilityを保つ。この構成により、同じ
versionの $P$ についてこれら全ての性質が得られた。

closed sourcesにも式~\eqref{eq:compensator-minimum}と同じminimum-stop compatibilityがあり、差を取ることで $(Q^n)$ にも
同じcompatibilityを得た。各 $Q^n$ は c\`adl\`ag true martingaleなので、martingale
貼り合わせ定理から、一つのstrongly adaptedな c\`adl\`ag local martingale $\widetilde Q$ を得る。この
$\widetilde Q$ は全ての $n$ について、$\tau_n$ で停止すると $Q^n$ とindistinguishableである。各
$\tau_n$ で
\[
 J^{\tau_n}\sim \widetilde Q^{\tau_n}+P^{\tau_n}
\]
が成り立つ。可算個のstopped agreementとexhaustionの共通full-measure set上では、任意の時刻 $t$ に
対し $t\leq\tau_n$ となる $n$ を選べる。従って $J\sim\widetilde Q+P$ である。正規化後の同じ
representative上で $Q=J-P$ と定めれば $Q\sim\widetilde Q$ である。従って $Q$ は
局所 martingaleであると同時に、適合 c\`adl\`ag 局所有限変動過程であり、$Q_0=0$ である。
この残差を局所 martingaleと示す際には、local propertyの時刻零indicatorを直接展開せず、
true-martingale coordinatesの貼り合わせとindistinguishabilityによる移送を用いる。

$L=N-Q=R+P$ と置けば $N=L+Q$ であり、$L$ も c\`adl\`ag 局所 martingaleである。しかし
$|\Delta R|\leq c$ に加え、次の補償過程の評価を用いる。

\begin{lemma}[補償過程の jump estimate]
\label{lem:compensator-jump-bound}
固定した $c>0$ と $T<\infty$ に対し、一つの full-measure set上で
\[
 0<t\leq T\quad\Longrightarrow\quad
 |\Delta P_t|\leq c,
 \qquad |\Delta L_t|\leq 2c .
\]
\end{lemma}

\begin{proof}
$R=L+(-P)$ を時刻 $T$ で停止する。予測可能 c\`adl\`ag 有限変動過程 $P^T$ の全ての非零jumpは、
累積変動の有理level crossingから作る可算個のannounced predictable stopping graphs
$(\sigma_k)$ で覆われる。完成した各時刻は $T+1$ 以下である。局所化と条件付き期待値の評価を
$L^T$ と $-P^T$ に適用すると、元の局所martingaleのjump全体の可積分性を仮定せずに
\[
 |\Delta (P^T)_{\sigma_k}|
 \leq
 \mathbb E[c\mid\mathcal F_{\sigma_k-}]=c
 \quad\text{a.s.}
\]
を得る。入力の評価は $|\Delta(R^T)_t|\leq c$ であり、$t>T$ では停止済み過程のjumpが零である。
可算個の零集合をまとめ、$(\sigma_k)$ のcoverageを用いれば、$P$ のjumpが零である時刻も含めて
最初の結論が従う。最後に $\Delta L=\Delta R+\Delta P$ と三角不等式を用いる。
$T=0$ は空な時間区間として分け、enumerationのdummy value $T+1$ では停止過程のjumpが零であることを使う。
\end{proof}

従って補償後に必要なestimateは、$Q=J-P$ 自身に対する $c$ boundではなく、まず
$|\Delta P|\leq c$、次にbounded-jump martingale $L=N-Q$ に対する $2c$ boundである。
これにより、元の零始点c\`adl\`ag local martingale $N$、$c>0$、有限horizon $T$、usual conditionsから
大jumpの局所化列と補償過程を構成し、同じ $L,Q$ に対する $N=L+Q$ を得る。
両成分は零始点c\`adl\`ag local martingaleであり、$Q$ は適合局所有限変動過程である。
$Q$ の予測可能性は主張しない。任意のpathwise stop $\rho\leq T$ についても、同じ零集合の外で
$|\Delta(L^\rho)_t|\leq2c$ が全時刻で成り立つ。これはjump評価の停止安定性である。
得られる分解は固定horizonに依存する。

\begin{lemma}[Horizon間の整合性]
各horizon $T$ で得た成分を $P^{(T)},Q^{(T)},L^{(T)}=N-Q^{(T)}$ と書く。
$U\leq T$ なら
\[
 (J^{c,T})^U=J^{c,U}
\]
が経路ごとに成り立ち、一つのfull-measure set上で全時刻について
\[
 (P^{(T)})^U=(P^{(U)})^U,\qquad
 (Q^{(T)})^U=(Q^{(U)})^U,\qquad
 (L^{(T)})^U=(L^{(U)})^U
\]
が成り立つ。各horizon内部のlocalizing sequenceを共通に選ぶ必要はない。
\end{lemma}

\begin{proof}
$(0,U]$ の大jump時刻集合は $(0,T]$ のものの部分集合である。
後者にのみ属する時刻は $U$ より大きく、時刻 $t\wedge U$ での有限和には寄与しない。
これにより最初の等式が従う。残りについては
\[
 A=(P^{(T)})^U-(P^{(U)})^U
   =(Q^{(U)})^U-(Q^{(T)})^U
\]
を考える。左辺は予測可能で、有限horizonで停止したため全時間軸で有限変動を持つ。
右辺はc\`adl\`ag local martingaleの差である。従ってpredictable FV local-martingale rigidityから
$A$ は初期値とindistinguishableである。初期値は零であるため、$P$ の一致が従い、同じ式から $Q$、
さらに $L=N-Q$ の一致を得る。$U=0$ は零初期値から直接従う。
\end{proof}

元の零始点local martingaleから各 $T_n=n+1$ の分解を構成し、上の補題をその同じ成分へ適用する。
可算交叉により、全 $m\leq n$ と全時刻でのoverlapを一つのfull-measure set上に揃えられる。
また同じ $P^{(T_n)}$ のjump boundも全 $n$ で共通のfull-measure set上に揃えられる。

\begin{lemma}[予測可能成分の全時刻への貼り合わせ]
元の零始点local martingaleから上記のcoherent familyと、
零始点の予測可能な右連続局所有限変動過程 $P$ を構成でき、
一つのfull-measure set上で全 $n,t$ に対して
$P_{t\wedge T_n}=P^{(T_n)}_{t\wedge T_n}$ が成り立つ。
さらに一つのfull-measure set上で全 $t$ に対し $|\Delta P_t|\leq c$ である。
\end{lemma}

\begin{proof}
各 $P^{(T_n)}$ を $T_n$ で停止して全時間軸で有限変動の座標とする。
二つのindexの大小に応じてoverlapを適用すれば、両座標は $T_n\wedge T_m$ で停止した後に一致する。
決定論的停止時刻列 $T_n$ は単調で無限大へ発散するため、予測可能な有限変動過程の貼り合わせを
適用できる。得られた過程の初期値はほとんど確実に零である。その例外集合はusual conditionsにより
時刻零で可測なので、その集合上で全pathを零に修正する。
可算交叉で全horizonの一致を同時に確保する。
任意の $t$ に対し $t\leq T_n$ となる $n$ を取り、停止前および停止時刻のleft jump一致から
$|\Delta P_t|=|\Delta P^{(T_n)}_t|\leq c$ を得る。
\end{proof}

\begin{lemma}[Local-martingale座標の共通細分]
零始点c\`adl\`ag local martingales $X^n$ が、ほとんど確実に単調でexhaustiveな停止時刻列 $(\tau_n)$ に対し、
$(X^n)^{\tau_n\wedge\tau_m}\sim(X^m)^{\tau_n\wedge\tau_m}$ を満たすとする。
このときc\`adl\`ag local martingale $Y$ が存在し、全 $n$ で
$Y^{\tau_n}\sim(X^n)^{\tau_n}$ である。
\end{lemma}

\begin{proof}
$X^n$ のlocalizing sequenceを $(\eta^n_j)_j$ とする。
可算な列の抽出により、$\sigma_n=\tau_n\wedge\eta^n_{k(n)}$ がほとんど確実に無限大へ発散するように
$k(n)$ を選べる。$\rho_n=\inf_{j\geq n}\sigma_j$ と置くと、filtrationの右連続性により停止時刻であり、
単調に無限大へ発散し、$\rho_n\leq\tau_n\wedge\eta^n_{k(n)}$ である。
零初期値はlocalizing stopに付く時刻零のindicatorを除去するため、optional stoppingから
$(X^n)^{\rho_n}$ はtrue martingaleである。
元のcompatible familyの点ごとの極限を $W$ とすると、$W^{\rho_n}\sim(X^n)^{\rho_n}$ である。
これにより細分した座標もcompatibleであり、true-martingale gluingを適用できる。
得られた $Y$ は全 $\rho_n$ で $W$ と一致するため、exhaustionから $Y\sim W$ である。
元の $\tau_n$ で停止し直すと所望の一致を得る。
\end{proof}

\begin{proposition}[全時刻のDoléans--Dade--Yen分解]
零始点c\`adl\`ag local martingale $N$ と $c>0$ に対し、零始点c\`adl\`ag local martingales
$L,Q$ が存在し、$N=L+Q$、$Q$ は局所有限変動であり、一つのfull-measure set上で全時刻について
$|\Delta L_t|\leq2c$ が成り立つ。
\end{proposition}

\begin{proof}
同じhorizon familyの $Q^{(T_n)}$ に共通細分とgluingを適用する。
各horizonでの一致からlocal BVと初期値零がほとんど確実に従う。
例外集合は時刻零で可測であり、その上でpathを零に修正して全pathのlocal BVと零初期値を得る。
同じfamilyから貼った $P,Q$ は、全horizonの共通零集合上で
$(J^{c,T_n})^{T_n}=(Q+P)^{T_n}$ を満たす。
$L=N-Q$ と置けば、$t\leq T_n$ に対して
$\Delta L_t=\Delta(N-J^{c,T_n})_t+\Delta P_t$ である。
任意の正の $t$ を含むhorizonを選び、それぞれのbound $c$ を加える。
時刻零のjumpは零である。$L$ の残りの性質はlocal martingalesの差から従う。
\end{proof}

この同じ大域 $L$ に次のquadratic constructionを適用する。
有限horizonの $L^T$ への適用も同じ構成の別の入力である。

\begin{lemma}[Bounded-jump local martingaleのquadratic局所化]
零始点c\`adl\`ag local martingale $M$ が、各整数horizon $n+1$ で一つのfull-measure set上の
jump bound $|\Delta M_t|\leq b_n$ を持つとする。ここで $b_n\geq0$ は定数である。
$M$ のlocal-martingale localizing sequenceを $(\eta_n)$ とし、
\[
 \theta_n=\inf\{t\geq0:|M_t|>n+1\},\qquad
 \rho_n=\eta_n\wedge\theta_n\wedge(n+1)
\]
と置く。この列はexhaustiveであり、各 $M^{\rho_n}$ はtrue martingaleで、全時刻共通に
\[
 |M^{\rho_n}_t|\leq n+1+b_n\quad\text{a.s.}
\]
が成り立つ。従って各座標の終端値は $L^2$ に属する。
\end{lemma}

\begin{proof}
absolute passageをhorizonで切った列と $(\eta_n)$ はともにlocalizing sequenceであり、その最小値も
同じ性質を持つ。零初期値により、$\eta_n=0$ のpathでもindicator付きlocalizing stopと通常のclosed stopが
一致する。従って $M^{\eta_n}$ はtrue martingaleである。$\rho_n\leq\eta_n$ でのbounded optional stoppingから
$M^{\rho_n}$ のtrue-martingale性を得る。passage直前の絶対値は $n+1$ 以下であり、境界で加わるjumpは
$b_n$ 以下であるため、表示した全時刻評価が従う。定数の $L^2$ majorantにより終端値の二乗可積分性を得る。
\end{proof}

各座標に有限horizon quadratic constructionを適用し、minimum-stop compatibilityとgluingを用いて
一つのglobal optional quadratic variation $V$ を得る。同じ $V$ のlocalizer rootと座標のterminal rootは
a.e. に一致し、その期待値costも一致する。各座標でDavis評価を適用すると
\[
 \mathbb E\sup_{t\leq n+1}|M^{\rho_n}_t|
 \leq 6\,\mathbb E\sqrt{V_{\rho_n}}
\]
を得る。ここでsupremumはc\`adl\`ag性により可算な時刻集合から検出する。

固定horizon DDY分解では $M=L^T$、$b_n=2c$ としてこの全構成を適用する。
元の零始点local martingaleからDDY成分を選び、その同じ $L^T$ に対するquadratic/Davis評価まで得られる。
この $V$ は $L^T$ のquadratic variationであり、元の $N$ のものではない。また各座標の評価から
unstopped $L^T$ のglobal $H^1$ 可積分性は主張しない。
